\documentclass[12pt]{article}

% ── Packages ─────────────────────────────────────────────────────────────────
\usepackage[T1]{fontenc}
\usepackage[utf8]{inputenc}
\usepackage{amsmath,amssymb,amsthm}
\usepackage{geometry}
\usepackage{hyperref}
\usepackage{booktabs}
\usepackage{enumitem}
\usepackage{algorithm}
\usepackage{algpseudocode}
\usepackage{microtype}
\usepackage{cite}

\geometry{margin=1in}

% ── Theorem environments ──────────────────────────────────────────────────────
\newtheorem{theorem}{Theorem}
\newtheorem{lemma}[theorem]{Lemma}
\newtheorem{corollary}[theorem]{Corollary}
\newtheorem{definition}[theorem]{Definition}
\newtheorem{proposition}[theorem]{Proposition}
\newtheorem{invariant}[theorem]{Invariant}

% ── Macros ───────────────────────────────────────────────────────────────────
\newcommand{\PHI}{\varphi}
\newcommand{\synBind}{\texttt{synBind}}
\newcommand{\synQuery}{\texttt{synQuery}}
\newcommand{\synRevoke}{\texttt{synRevoke}}

% ─────────────────────────────────────────────────────────────────────────────
\title{\textbf{Self-Healing Multi-Agent Systems via \\
  Permanent State Imprinting and Autonomous Job Queues}}

\author{Alfredo Medina Hernandez \\
  \normalsize Medina Tech \quad Dallas, Texas, USA \\
  \normalsize \texttt{MedinaSITech@outlook.com}}

\date{April 2026}
% ─────────────────────────────────────────────────────────────────────────────

\begin{document}

\maketitle

\begin{abstract}
We present the \emph{Synapse Binding Engine} (SYN), a mechanism for eliminating
cross-agent query latency in multi-agent systems through a \emph{permanent
contract}: a single cross-canister call imprints a remote agent's state
snapshot into the local canister's stable memory, after which all queries are
pure local reads with zero network cost.  We formalise three properties of SYN
— \emph{imprint permanence}, \emph{upgrade survival}, and \emph{owner-revocable
termination} — and prove that a SYN-bound system achieves eventual consistency
with bounded staleness proportional to the GOVERNANCE\_SYNC job interval.  We
then describe the autonomous job queue that drives self-healing: five job types
($\mathsf{ADD\_HOTKEY}$, $\mathsf{AUTO\_DISCOVER}$, $\mathsf{GOVERNANCE\_SYNC}$,
$\mathsf{DEPLOY\_WORKER}$, $\mathsf{HEARTBEAT\_CHECK}$) with four priority
levels and exponential back-off retry, governed by live HEART-snapshot routing
decisions.  We show that the resulting system is \emph{self-bootstrapping},
\emph{self-healing}, and \emph{self-governing} without any human intervention
after a single master-boot call.
\end{abstract}

% ─────────────────────────────────────────────────────────────────────────────
\section{Introduction}
\label{sec:intro}

The dominant paradigm for cross-agent coordination is \emph{on-demand querying}:
whenever Agent~$A$ needs data from Agent~$B$, it sends a message, waits for the
reply, and then proceeds.  This paradigm has three costs that compound in
high-frequency autonomous systems:

\begin{enumerate}[label=(\arabic*)]
  \item \textbf{Latency.}  Even on high-throughput blockchains, a cross-canister
    call introduces a round-trip delay of hundreds of milliseconds.  In a system
    ticking at 875\,ms, one avoidable cross-canister call consumes
    $\approx\!30\%$ of the tick budget.
  \item \textbf{Cycle cost.}  Every cross-canister call consumes compute units.
    A system that queries its ``brain'' 50 times per tick at 43-second intervals
    spends more on coordination than on computation.
  \item \textbf{Coupling fragility.}  If Agent~$B$ is temporarily unreachable,
    Agent~$A$'s dependent computation fails — even if the relevant state has
    not changed.
\end{enumerate}

SYN addresses all three by converting the first cross-agent query into a
\emph{permanent local write}.  After the imprint, queries cost nothing; the
state survives canister upgrades; and temporary unreachability of the remote
agent does not affect the local agent's operation.

Section~\ref{sec:model} establishes the system model.
Section~\ref{sec:syn} defines SYN and proves its key properties.
Section~\ref{sec:jobqueue} describes the autonomous job queue.
Section~\ref{sec:selfhealing} analyses the self-healing dynamics.
Section~\ref{sec:impl} reports implementation details.
Section~\ref{sec:related} situates the work relative to existing distributed
systems literature.

% ─────────────────────────────────────────────────────────────────────────────
\section{System Model}
\label{sec:model}

\begin{definition}[Canister]
  A \emph{canister} is an autonomous actor with: (i) stable state $\Sigma$ that
  persists across upgrades via a serialisation/deserialisation hook;
  (ii) a message queue from which it processes calls sequentially;
  (iii) an owner principal $p_\text{own}$ that controls administrative
  operations; and (iv) a globally unique canister ID $\mathrm{id}$.
\end{definition}

\begin{definition}[HEART Snapshot]
  The \emph{HEART snapshot} $H = (b, \tau, \eta, s, m, v, \pi, o, c, \sigma)$
  is a 10-tuple of the AGI Terminal canister's key state variables:
  \begin{align*}
    b &\in \{0,1\} && \text{booted flag} \\
    \tau &\in \mathbb{N} && \text{current tick} \\
    \eta &\in \mathbb{N} && \text{neuron count} \\
    s &\in \mathbb{N} && \text{total staked ICP (e8s)} \\
    m &\in \mathbb{N} && \text{pending maturity (e8s)} \\
    v &\in \mathbb{R}_{\ge 0} && \text{voting power} \\
    \pi &\in \mathbb{N} && \text{treasury ICP (e8s)} \\
    o &\in \mathbb{N} && \text{ONESICAN inventory} \\
    c &\in \mathbb{N} && \text{circulating OnesICANs} \\
    \sigma &\in \mathbb{N} && \text{spawned neurons}
  \end{align*}
\end{definition}

\begin{definition}[SYN Binding]
  A \emph{SYN binding} is a tuple
  $B = (\ell, \mathrm{cid}, k, d, t_b, n_r)$ stored in the local canister's
  stable memory, where:
  $\ell$ is the binding label (a string key);
  $\mathrm{cid}$ is the remote canister ID;
  $k$ is the data key passed to the remote canister;
  $d$ is the locally stored data snapshot (serialised text);
  $t_b$ is the timestamp of the last bind call;
  $n_r \in \mathbb{N}_0$ is the refresh count.
\end{definition}

% ─────────────────────────────────────────────────────────────────────────────
\section{Synapse Binding Engine}
\label{sec:syn}

\subsection{Operations}

SYN exposes five public operations:

\begin{description}[leftmargin=1.8cm, labelwidth=1.6cm]
  \item[\synBind$(\ell, \mathrm{cid}, k)$]
    Performs one cross-canister call to $\mathrm{cid}$ to retrieve the current
    value for key $k$; stores the result as a binding with label $\ell$ in
    local stable memory.  All future reads use the local copy.
  \item[\synQuery$(\ell)$]
    Returns the locally stored data for label $\ell$.
    \textbf{No network call.  No cycle cost.  Deterministic.}
  \item[\synRevoke$(\ell)$]
    Owner-only.  Deletes the binding for $\ell$.  Data is gone; no recovery.
  \item[$\synRevoke$All$()$]
    Owner-only.  Deletes all bindings.  Solver goes dark on the next tick.
  \item[\texttt{synBindHeart}$(\mathrm{cid}_\text{AGI})$]
    Specialisation of \synBind{} for the HEART snapshot: calls
    $\mathrm{cid}_\text{AGI}.\texttt{getSystemStatus}()$ and imprints all ten
    fields locally.  Also caches $m$ (maturity) and $v$ (voting power) as
    first-class numeric fields for fast routing decisions without string parsing.
\end{description}

\subsection{Invariants}

\begin{invariant}[Imprint Permanence]
  After \synBind$(\ell, \mathrm{cid}, k)$ returns successfully, $\synQuery(\ell)$
  returns the imprinted data without any network call, for the lifetime of the
  canister (unless \synRevoke\ is called by the owner).
\end{invariant}

\begin{invariant}[Upgrade Survival]
  Because bindings are stored in the canister's stable memory and serialised
  through the pre-upgrade/post-upgrade hooks, all bindings survive an arbitrary
  number of canister upgrades.
\end{invariant}

\begin{invariant}[Owner-Gated Termination]
  Only the owner principal $p_\text{own}$ may call \synRevoke\ or
  $\synRevoke$All.  No binding can be deleted by any other principal, including
  the remote canister from which the data was originally retrieved.
\end{invariant}

\subsection{Staleness Analysis}

\begin{theorem}[Bounded Staleness]
  Under the PRPO governance-sync schedule with interval $\Delta_\text{GS}$
  (seconds), the staleness of any SYN binding relative to the current remote
  state is bounded by $\Delta_\text{GS}$ in expectation.  Specifically, if the
  remote state changes at rate $\lambda$ (events/second), then the expected
  number of unseen updates in a binding is $\lambda \cdot \Delta_\text{GS}$.
\end{theorem}

\begin{proof}
  Let $T_\text{bind}$ be the last bind time.  The binding reflects remote state
  at $T_\text{bind}$.  The next GOVERNANCE\_SYNC job refreshes the binding at
  $T_\text{bind} + \Delta_\text{GS}$ (in the common case; see the retry
  analysis in Section~\ref{subsec:retry} for the failure case).  In the
  interval $[T_\text{bind}, T_\text{bind} + \Delta_\text{GS}]$, the remote
  state undergoes $\lambda \cdot \Delta_\text{GS}$ expected changes.
  Therefore, the expected staleness (unseen updates) is $\lambda \Delta_\text{GS}$.
  \qed
\end{proof}

For the HEART binding, $\Delta_\text{GS}$ is set by the GOVERNANCE\_SYNC job
cadence; at one job per solver tick (50 AGI heartbeats, $\approx 43\,\text{s}$),
$\Delta_\text{GS} = 43\,\text{s}$.  For economic state variables (maturity,
treasury), $\lambda \approx 1/2\,\text{s}^{-1}$ (one ICP heartbeat per 2\,s),
giving expected staleness $\approx 21$ unobserved ticks per sync period.

\subsection{Capacity}

The Organism Solver supports up to $N_{\max} = 64$ simultaneous bindings,
limited by the stable memory budget for the binding array.  This is sufficient
to mirror the entire 39-canister \textsc{Nova} organism plus 25 future
expansion canisters.

% ─────────────────────────────────────────────────────────────────────────────
\section{Autonomous Job Queue}
\label{sec:jobqueue}

\subsection{Job Types and Semantics}

Five job types implement the full self-management lifecycle:

\begin{table}[h]
  \centering
  \caption{Job types, default priorities, and semantics.}
  \label{tab:jobs}
  \begin{tabular}{llp{6cm}}
    \toprule
    Job type & Default priority & Semantic \\
    \midrule
    \texttt{AUTO\_DISCOVER}   & CRITICAL & Map the current canister topology \\
    \texttt{GOVERNANCE\_SYNC} & HIGH     & Refresh fleet $+$ AI SYN bindings \\
    \texttt{HEARTBEAT\_CHECK} & NORMAL   & Ping all registered canisters \\
    \texttt{ADD\_HOTKEY}      & NORMAL   & Register a new NNS hotkey \\
    \texttt{DEPLOY\_WORKER}   & NORMAL   & Spawn a named service worker \\
    \bottomrule
  \end{tabular}
\end{table}

\subsection{Priority Scheduler}

Jobs are partitioned by priority (CRITICAL, HIGH, NORMAL, LOW) and sorted
within each partition by creation time.  At each solver tick, up to
$N_\text{tick} = 5$ jobs are processed; CRITICAL jobs bypass the limit.
When the queue reaches $N_\text{max} = 200$, LOW-priority jobs are evicted
first; if still full, NORMAL jobs are evicted.  CRITICAL and HIGH jobs are
never dropped.

\begin{algorithm}[h]
  \caption{Solver tick scheduler}
  \label{alg:tick}
  \begin{algorithmic}[1]
    \State $\mathit{eligible} \gets \{j \in Q : j.\mathit{status} \in \{\mathsf{PENDING}, \mathsf{RETRYING}\}
                   \wedge j.\mathit{nextRunAt} \le t_\text{now}\}$
    \State $\mathit{sorted} \gets \text{sort by priority}(\mathit{eligible})$
    \State $n \gets 0$
    \For{$j \in \mathit{sorted}$}
      \If{$n < N_\text{tick}$ \textbf{or} $j.\mathit{priority} = \mathsf{CRITICAL}$}
        \State $\text{execute}(j)$; $n \gets n + 1$
        \If{failed and $j.\mathit{retries} < R_\text{max}$}
          \State $j.\mathit{nextRunAt} \gets t_\text{now} + 2^{j.\mathit{retries}}\,\text{s}$
          \State $j.\mathit{status} \gets \mathsf{RETRYING}$
        \ElsIf{failed}
          \State $j.\mathit{status} \gets \mathsf{FAILED}$
          \If{$j.\mathit{type} = \mathsf{HEARTBEAT\_CHECK}$}
            \State \textbf{enqueue} $\mathsf{HEARTBEAT\_CHECK}$ at $\mathsf{CRITICAL}$
          \EndIf
        \EndIf
      \EndIf
    \EndFor
  \end{algorithmic}
\end{algorithm}

\subsection{Exponential Back-off Retry}
\label{subsec:retry}

Failed jobs retry with delay $2^{r}\,\text{s}$ where $r \in \{1, 2, 3\}$ is
the retry attempt.  At $R_\text{max} = 3$ failures, the job transitions to
$\mathsf{FAILED}$.  For $\mathsf{HEARTBEAT\_CHECK}$ failures — indicating a
possibly dead canister — a new $\mathsf{HEARTBEAT\_CHECK}$ at CRITICAL priority
is immediately re-queued.

\begin{proposition}[Retry Convergence]
  In the absence of permanent failure (i.e.\ if the remote canister recovers
  within $2^{R_\text{max}}\,\text{s} = 8\,\text{s}$), every job eventually
  completes.  In the presence of permanent failure, the job transitions to
  $\mathsf{FAILED}$ after at most $1 + 2 + 4 = 7$ seconds of cumulative retry
  delay.
\end{proposition}

% ─────────────────────────────────────────────────────────────────────────────
\section{Self-Healing Dynamics}
\label{sec:selfhealing}

\subsection{HEART-Driven Routing Decisions}

The HEART binding's cached maturity field $m$ drives autonomous priority
escalation at each solver tick:

\[
  \text{If } m > m_\text{threshold}:
  \quad \forall j \in Q,\; j.\mathit{type} \in \{\mathsf{ADD\_HOTKEY},
  \mathsf{GOVERNANCE\_SYNC}\}:
  \quad j.\mathit{priority} \gets \mathsf{CRITICAL}.
\]

At $m_\text{threshold} = 10^{10}\,\text{e8s}$ (100\,ICP worth of maturity),
the system recognises that economic activity has accumulated and upgrades
governance-related jobs to maximum priority.  This requires no human
observation: the HEART snapshot, refreshed every 43\,s, provides the signal;
the solver acts on it without human input.

\subsection{Self-Bootstrapping via \texttt{initialize()}}

The \texttt{initialize()} call queues the three bootstrap jobs in a single
transaction:
$\mathsf{AUTO\_DISCOVER}$ (CRITICAL),
$\mathsf{GOVERNANCE\_SYNC}$ (HIGH),
$\mathsf{HEARTBEAT\_CHECK}$ (NORMAL).
After this call, the system requires no further configuration: the AGI Terminal
fires $\texttt{solverTick()}$ every 50 system ticks, the jobs execute, SYN
bindings are populated, and the organism is operational.

\subsection{Canister Self-Diagnosis via Syntax Synapse}

A companion canister, the Syntax Synapse Engine, maintains a registry of all 39
canisters and periodically audits them for error patterns.  It classifies six
error categories:

\begin{table}[h]
  \centering
  \caption{Syntax Synapse error pattern taxonomy.}
  \label{tab:synapse}
  \begin{tabular}{ll}
    \toprule
    Pattern tag & Source \\
    \midrule
    \texttt{\#reservedKeyword}    & Use of Motoko reserved word as field name \\
    \texttt{\#paramSeparator}     & Semicolon instead of comma in param list \\
    \texttt{\#piConstant}         & Float literal without decimal point \\
    \texttt{\#missingParens}      & Unwrapped switch/if condition \\
    \texttt{\#intFromNat}         & Type mismatch between \texttt{Int} and \texttt{Nat} \\
    \texttt{\#lambdaType}         & Lambda type without enclosing parentheses \\
    \bottomrule
  \end{tabular}
\end{table}

For each error pattern, the Syntax Synapse Engine returns a fix recommendation
via \texttt{recommendFix(pattern)}.  The system heartbeat (every 13 ticks)
re-audits the registry and updates status to ``OK'' or ``ERROR''.  This closes
a second self-healing loop independent of SYN.

% ─────────────────────────────────────────────────────────────────────────────
\section{Implementation Details}
\label{sec:impl}

\subsection{Stable Memory Layout}

The Organism Solver uses the \texttt{persistent actor} pattern from the Motoko
base library, which automatically serialises all \texttt{var} declarations
(equivalent to \texttt{stable var}) through the pre-upgrade/post-upgrade hooks.
The binding array \texttt{\_bindings : [SynBinding]} occupies a contiguous
region of stable memory that survives an unlimited number of canister upgrades.

\subsection{Re-entrancy Guard}

The \texttt{solverTick()} function is guarded by a transient boolean
\texttt{\_tickRunning}.  Because the guard is \texttt{transient} (not stable),
it resets to \texttt{false} after every upgrade — ensuring that a crashed tick
cannot permanently lock the solver.

\subsection{Ring Buffer Event Log}

The solver maintains a ring buffer of the last $L = 100$ events.  When the
log exceeds $L$ entries, the oldest entries are discarded:
\[
  \mathtt{\_log}[t] = \mathtt{\_log}[(t \bmod L)].
\]
This provides a lightweight audit trail without unbounded memory growth.

\subsection{Complexity}

$\synBind$ executes in $\order{1}$ local operations plus one cross-canister
round trip.  $\synQuery$ executes in $\order{|\mathtt{\_bindings}|}$ time
(linear scan for the label), bounded by $N_\text{max} = 64$.
$\synRevoke$ is also $\order{N_\text{max}}$.  The scheduler tick is
$\order{Q_\text{eligible} \log Q_\text{eligible}}$ for the sort, then $\order{N_\text{tick}}$
for execution.

% ─────────────────────────────────────────────────────────────────────────────
\section{Related Work}
\label{sec:related}

\paragraph{Actor model.}
Hewitt's actor model~\cite{hewitt1973} establishes message-passing as the
primitive for concurrent computation.  SYN complements this by providing a
mechanism to \emph{eliminate} message-passing for stable read-heavy data.

\paragraph{Eventual consistency.}
Vogels~\cite{vogels2009} surveys eventual consistency in large-scale distributed
systems.  SYN achieves a weaker form — \emph{sync-periodic consistency} — where
staleness is bounded by the governance-sync interval rather than propagation
delay.

\paragraph{Content-addressed storage.}
IPFS~\cite{benet2014} achieves immutability through content addressing.  SYN
achieves a different form of immutability: the binding data is mutable (can be
refreshed) but the binding's existence is permanent until owner revocation.

\paragraph{Distributed caches.}
Redis and Memcached provide distributed caching with configurable expiry.  SYN
differs in three ways: (i) bindings have no automatic expiry; (ii) the binding
survives canister upgrades; (iii) only the owner can delete a binding.

% ─────────────────────────────────────────────────────────────────────────────
\section{Conclusion}
\label{sec:conclusion}

The Synapse Binding Engine eliminates the fundamental cost of cross-agent
coordination: the recurring query.  By converting the first query into a
permanent local write, SYN enables a multi-agent system to operate with bounded
staleness, zero coordination latency for all reads after the initial bind, and
owner-controlled termination semantics that do not depend on the availability
of the remote agent.

The autonomous job queue ensures that the binding refresh cycle, the canister
health check cycle, and the governance synchronisation cycle all operate without
human input.  The HEART snapshot closes the feedback loop: the system reads its
own economic state and adjusts its job priorities accordingly.  The result is a
self-bootstrapping, self-healing, self-governing organism that requires a single
master-boot call and then runs forever.

% ─────────────────────────────────────────────────────────────────────────────
\begin{thebibliography}{99}

\bibitem{hewitt1973}
C.~Hewitt, P.~Bishop, R.~Steiger,
``A universal modular actor formalism for artificial intelligence,''
\textit{Proc.\ 3rd IJCAI}, pp.~235--245, 1973.

\bibitem{vogels2009}
W.~Vogels,
``Eventually consistent,''
\textit{Queue}, vol.~6, no.~6, pp.~14--19, 2009.

\bibitem{benet2014}
J.~Benet,
``IPFS -- Content addressed, versioned, P2P file system,''
arXiv:1407.3561, 2014.

\bibitem{armstrong2010}
J.~Armstrong,
\emph{Programming Erlang: Software for a Concurrent World},
2nd ed., Pragmatic Bookshelf, 2013.

\bibitem{lamport1978}
L.~Lamport,
``Time, clocks, and the ordering of events in a distributed system,''
\textit{Commun.\ ACM}, vol.~21, no.~7, pp.~558--565, 1978.

\bibitem{dfinity2022}
\text{DFINITY Foundation},
``The Internet Computer,''
Technical White Paper, 2022.

\bibitem{gilbert2002}
S.~Gilbert, N.~Lynch,
``Brewer's conjecture and the feasibility of consistent, available,
partition-tolerant web services,''
\textit{SIGACT News}, vol.~33, no.~2, pp.~51--59, 2002.

\end{thebibliography}

\end{document}
